\section{Introduction}
\label{sec:introduction}

\begin{figure}[t]
\noindent\fbox{\parbox{0.97\linewidth}{\small
\textbf{Summary for student-support staff.}
\emph{The decision.} After the first semester, a support office can meet only about one in five
students and must decide whom to contact first.
\emph{What the model does.} It assigns each student a dropout-risk score from records the
university already holds (first-semester results, fees and background) and places the top 20\% on
a contact list.
\emph{Main findings.} A simpler, fully transparent model performed as well as the complex one. On
past data, the list was highly precise: 108 of the 109 flagged test students dropped out. Yet most
flagged validation students had passed no units, which staff can already see, and the list misses
about half of the students who later leave, largely because only one in five can be contacted. The
displayed ``reasons'' are useful hints, but they may change when recalculated, and there was no
evidence that seeing them helped people anticipate the model's decision.
\emph{Limitations.} The model was tested only on Portuguese data from students with a known final
outcome. It does not explain \emph{why} a student is struggling or what will help, and it must
never be used to deny support.}}
\end{figure}

\subsection{Background and Problem Statement}

Student dropout is costly for students and institutions alike, and the dataset used here was
created to help identify at-risk students early enough for support~\citep{realinho2021uci,
realinho2022predicting}. In practice, a student-support office can hold only a limited number of
outreach meetings after the first semester and must therefore rank students and contact those at
highest risk first. A ranking alone is insufficient, however: before a conversation, an adviser also
needs to understand \emph{why} a student was flagged, and those reasons must be trustworthy.

Tabular foundation models such as TabPFN~\citep{hollmann2025tabpfn} and
\tabicl~\citep{qu2025tabicl} have been reported to match or outperform tuned gradient boosting on
many small tabular benchmarks without dataset-specific training. Yet they are black boxes whose
explanations must be obtained post hoc, for example with SHAP~\citep{lundberg2017shap}, and
\citet{rudin2019stop} argues that an interpretable model should be preferred for high-stakes
decisions whenever it performs comparably. Against this background, the present study treats
\tabicl not as a model to be promoted but as a test case for whether black-box complexity is
warranted in student outreach.

\textbf{Problem statement.} Benchmark accuracy does not establish that a foundation model improves a
capacity-limited outreach decision, nor that its post-hoc explanations are faithful, stable or
useful to the people who must act on them. None of the representative studies in
Table~\ref{tab:prior-work} combines these three perspectives for dropout outreach. The project
therefore asks whether a support office should use a black-box model for this decision and whether
its explanations withstand both technical and human evaluation.

\subsection{Decision Setting}

The intended user is a student-support officer rather than the student. Predictions are made once,
after the first semester; second-semester information is excluded, as it would not yet be
available. The operating rule was fixed in advance: all students in a cohort are ranked by
predicted dropout probability, and exactly the top 20\% (the assumed outreach capacity) are flagged,
with ties broken by record order. A flag means ``offer a support conversation''; it never implies
withdrawing support or making an automated decision about the student. Errors carry different
costs: a missed student may leave without being offered help, an unnecessary flag consumes a scarce
meeting and may stigmatise a student who was coping, and a misread explanation may steer the
conversation towards the wrong problem.

\subsection{Research Questions and Contributions}

\begin{enumerate}[nosep,leftmargin=*]
  \item[\textbf{RQ1.}] Does the black-box \tabicl add predictive value over simpler conventional
  and intrinsically interpretable baselines at a fixed 20\% outreach capacity?
  \item[\textbf{RQ2.}] Are its local explanations faithful, stable and appropriately constrained,
  and what do they reveal to an adviser about the students it flags?
  \item[\textbf{RQ3.}] Do explanation charts help people anticipate the model's outreach flag and
  identify useful support information?
\end{enumerate}

\noindent To answer these questions, the project makes four contributions:
\begin{enumerate}[nosep,leftmargin=*]
  \item a pre-specified, locked comparison of \tabicl with logistic regression, a depth-3 decision
  tree, XGBoost and a main-effects EBM at the same outreach capacity, with paired bootstrap
  intervals \textbf{(RQ1)};
  \item a post-hoc audit of the explanations for all 544 validation students, covering
  faithfulness, stability and constrained recourse, together with an analysis of which students are
  flagged \textbf{(RQ2)};
  \item a two-group forward-simulation study with 50 adult participants, including a qualitative
  analysis of 1,500 written answers \textbf{(RQ3)}; and
  \item sensitivity analyses for still-enrolled students, an admission-time feature window and a
  model without protected attributes, accompanied by a local, validation-only research interface.
\end{enumerate}

\textbf{Scope.} The study relies on one public dataset from a single Portuguese institution, a
binary Dropout-versus-Graduate target and one fixed data split, and its fairness analysis is
descriptive. The human study involved volunteers acting as proxy users; it thus measures
understanding of the model's output, not decisions about real students. Table~\ref{tab:terms}
defines the key terms used throughout the report.

\begin{table}[!htb]
\centering
\caption{Key terms used in this report.}
\label{tab:terms}
\footnotesize
\setlength{\tabcolsep}{3pt}
\begin{tabular}{@{}>{\raggedright\arraybackslash}p{2.0cm}p{6.0cm}@{}}
\toprule
\textbf{Term} & \textbf{Meaning} \\
\midrule
Outreach capacity & The proportion of students the university can contact or support; assumed in this study to be 20\% of a cohort \\
Flag (flagged student) & A student ranked by the model among the 20\% with the highest predicted dropout risk, who may need additional attention and would be offered a support meeting \\
Risk score & The model's estimated dropout probability; a score, not a certainty \\
Contribution (SHAP) & How far one group of information moved a student's score relative to a typical student; association, not cause \\
Faithfulness & Whether the reasons shown are those the model actually relies on \\
Stability & Whether the same top five reasons reappear when the explanation is recalculated \\
Recourse & Actions that a student or university staff could still take on changeable factors (in this study, only fee or debt status; completed first-semester results are fixed) that would lower the model's predicted risk or remove the flag; these show model behaviour, not proven ways to prevent dropout \\
Forward simulation & People predict the model's decision from the same information \\
Locked final test & Held-out data used exactly once, after all choices were fixed \\
\bottomrule
\end{tabular}
\end{table}
